% ---
% id: receiver/thm:emitted-bound
% kind: thm
% title: "Multiclick bounds carried by the emitted weights"
% summary: >-
%   Each setting's emitted weight bounds its multiclick element on every
%   photon-number block above the cutoff, and their minimum does so at every
%   setting, under the eigenvalue-slack assumption [[A-EIG]] only where
%   computed block eigenvalues enter. Each weight lies in [0, 1] and is
%   finitely computable.
% prominence: main
% status: implemented
% uses: [def:setting-family, def:pmc, def:pmcex, def:enclosed-weight, thm:enclosed-weight, thm:multiclick-bound, def:exactbnd, lem:collapse, lem:windowtail, lem:exactwindow, lem:exactbasic]
% mentions: [rem:exacteig, rem:exactpositivity, req:rigorous-weight]
% assumes: [A-EIG]
% refs: []
% provenance: newly-derived
% implements: certiqs_key.receiver.bound:emitted_weights
% lean:
% notes: "The per-setting inequality eq:pmcsetting is the first conclusion and the uniform eq:pmcstatement its consequence; a claim that cites eq:pmcstatement reads the uniform form under its own label."
% ---

This is the theorem the emitted weights rest on. Each setting's weight is
the analytic weight of Definition~\ref{def:pmc} or, where a depth is
declared, the exact-window weight of Definition~\ref{def:pmcex}; where the
engine evaluates the weight rigorously (Requirement~\ref{req:rigorous-weight})
it is the enclosed weight of Definition~\ref{def:enclosed-weight}, formed by
the same formulas from bounds on the conservative side. The theorem says that
each bounds the multiclick element of its own setting on every
photon-number block above the cutoff, and that the minimum over the settings
bounds it at every setting, under Assumption~\ref{ass:eig} only where a
computed block eigenvalue enters, at a setting governed by
Definition~\ref{def:pmcex}. Remark~\ref{rem:exacteig} explains why that entry
rests on the assumption rather than on a theorem, and
Theorem~\ref{thm:enclosed-weight} is the theorem that replaces it.

\begin{theorem}\label{thm:emitted-bound}
Let a setting family (Definition~\ref{def:setting-family}) be given. For
every setting $t\in\Tset$ let $\pmc^{(t)}$ be as in
Definition~\ref{def:pmc} if $t$ carries no exact-window depth and as in
Definition~\ref{def:pmcex} if it carries one, or, at any setting, let it be
the \tref{enclosed weight} $\check p^{(t)}_{\mc}$ of
Definition~\ref{def:enclosed-weight} formed from \tref{enclosed bound data} of
$t$ and, where $t$ carries a depth, from \tref{rigorous window bounds} of $t$,
or a real $p$ with $0\le p\le\check p^{(t)}_{\mc}$;
and let $\pmc=\min_{t\in\Tset}\pmc^{(t)}$ as in Definition~\ref{def:pmc}. Then
\begin{equation}\label{eq:pmcsetting}
  \blk{M}{\gtN}_{\mc|t}\ \succeq\ \pmc^{(t)}\,\id^{(>N)}
  \qquad\text{for every }t\in\Tset ,
\end{equation}
where $\blk{M}{\gtN}_{\mc|t}$ is the restriction of the multiclick element
$M_{\mc|t}$ of \eqref{eq:family-elements} at setting $t$ to the photon-number
blocks $n>N$; and consequently, with the one weight $\pmc$ at every setting,
\begin{equation}\label{eq:pmcstatement}
  \blk{M}{\gtN}_{\mc|t}\ \succeq\ \pmc\,\id^{(>N)}
  \qquad\text{for every }t\in\Tset .
\end{equation}
At a setting governed by Definition~\ref{def:pmc} or by
Definition~\ref{def:enclosed-weight} both inequalities hold at
that setting with no hypothesis beyond this contract's ledger; at a setting
governed by Definition~\ref{def:pmcex} they hold there \emph{under
Assumption~\ref{ass:eig}}, which is what makes the second entry of
\eqref{eq:elldef} a lower bound. At a
one-setting family, $\Tset=\{\off\}$, one has $\pmc=\pmc^{(\off)}$ and the
two inequalities are one. Moreover $0\le\pmc^{(t)}\le1$ at every setting,
hence $0\le\pmc\le1$, and each $\pmc^{(t)}$ of Definition~\ref{def:pmc}
or~\ref{def:pmcex} is computable by finitely many
arithmetic operations on
$\{\lmax(B^{(t)}),\lmin(B^{(t)}_k),\gamma^{(t)},\beta^{(t)}_k\}$ and the
declared $N$, $m_t$ and $m^{\ast}_t$ --- where $B^{(t)},B^{(t)}_k$ are the
operators of the setting (Definition~\ref{def:setting-family}\,(b)) ---
together with, where the setting carries a declared depth, the declared
$\varepsilon_{\mathrm{eig}}$ and the $m^{\ast}_t-N$ extreme eigenvalues
$\widehat\Bnd^{\ast(t)}(n)$ of the Hermitian matrices \eqref{eq:exactbnd},
each of dimension $\binom{D+n-1}{n}$; each enclosed weight by finitely many
on the enclosed bound data, the rigorous window bounds and the declared $N$,
$m_t$ and $m^{\ast}_t$; and $\pmc$ by finitely many comparisons among them.
\end{theorem}

\begin{proof}
\emph{Idea.} At a fixed setting, Theorem~\ref{thm:multiclick-bound} bounds the
multiclick element block by block. Lemma~\ref{lem:collapse} or
Lemma~\ref{lem:windowtail} turns the block bounds into the analytic weight,
and Lemma~\ref{lem:exactwindow} turns them into the exact-window weight;
Theorem~\ref{thm:enclosed-weight} does both for the enclosed weight. The
uniform statement follows because a minimum lies below each of its terms.

\emph{Step 1 (the analytic weight at a fixed setting).} Fix $t$. By
Theorem~\ref{thm:multiclick-bound}, applicable at every setting by
\eqref{eq:family-mc}, $\blk{M}{\eqn}_{\mc|t}\succeq\Bnd^{(t)}(n)\id$ for
every $n$, so
$\blk{M}{\gtN}_{\mc|t}\succeq\big(\inf_{n>N}\Bnd^{(t)}(n)\big)\id^{(>N)}$
--- the blocks are orthogonal summands and a lower bound holding on each
holds on the direct sum. Since $\blk{M}{\gtN}_{\mc|t}\succeq0$ as well
($M_{\mc|t}$ is a sum of POVM elements,
Definition~\ref{def:setting-family}\,(a)), the bound may be replaced by its
positive part. In the first branch of \eqref{eq:pmcdef} --- where the flag
is $\mathrm{exact}$, so that \eqref{eq:collapsehyp} holds with equality by
construction (Definition~\ref{def:setting-family}\,(c)) ---
Lemma~\ref{lem:collapse} identifies that positive part with
$\max\{0,\Bnd^{(t)}(N+1)\}$; in the second,
Lemma~\ref{lem:windowtail}\,(i) gives
$\underline{\Bnd}^{(t)}_{N,m_t}\le\inf_{n>N}\Bnd^{(t)}(n)$, so its positive
part is a valid lower bound as well.

\emph{Step 2 (the exact-window weight at a fixed setting).} At a setting
carrying a declared depth, \eqref{eq:exactstatement} of
Lemma~\ref{lem:exactwindow} gives the conclusion directly for
\eqref{eq:pmcex}; its hypothesis \eqref{eq:mcpos} holds at every setting,
because $M_{\mc|t}=\sum_{S\in\MCK}M_{S|t}$ is a sum of elements of a POVM
(Definition~\ref{def:setting-family}\,(a)), which is the discharge
Remark~\ref{rem:exactpositivity} says is needed. On branch~1, apply it with
$m^{\ast}=m^{\ast}_t$, $\ell=\ell^{(t)}$ and
$\tau=\Bnd^{(t)}(m^{\ast}_t+1)$, the remainder being admissible by
part~(ii) because predicate~(X) makes \eqref{eq:collapsehyp} hold. On
branch~2, apply it with $m^{\ast}=m_t$, $\tau=T^{(t)}_{m_t}$ admissible by
part~(i), and the window function $\ell(n)=\ell^{(t)}(n)$ for
$N<n\le m^{\ast}_t$, $\ell(n)=\Bnd^{(t)}(n)$ for $m^{\ast}_t<n\le m_t$ ---
the two ranges together being $N<n\le m_t$ because $m^{\ast}_t\le m_t$ is
required there. In both cases hypothesis \eqref{eq:ellhyp} holds entry by
entry: for $\Bnd^{(t)}(n)$ by Lemma~\ref{lem:exactbasic}, and for
$\widehat\Bnd^{\ast(t)}(n)-\varepsilon_{\mathrm{eig}}$ by
Assumption~\ref{ass:eig}; a maximum of the two is admissible by
part~(iii).

\emph{Step 2a (the enclosed weight at a fixed setting).} At a setting whose
weight is the enclosed weight $\check p^{(t)}_{\mc}$ or a real $p$ with
$0\le p\le\check p^{(t)}_{\mc}$,
Theorem~\ref{thm:enclosed-weight}\,(iii) gives \eqref{eq:pmcsetting} at $t$
and $0\le\pmc^{(t)}\le1$, for the enclosed weight and for every real between
$0$ and it, with no hypothesis beyond Definition~\ref{def:setting-family},
whether or not $t$ carries a depth.

\emph{Step 3 (the uniform weight).} So \eqref{eq:pmcsetting} holds at $t$,
under Assumption~\ref{ass:eig} exactly where $t$ is governed by
Definition~\ref{def:pmcex}. Since
$\pmc$ is the minimum of the $\pmc^{(s)}$ over the finite set $\Tset$,
$\pmc\le\pmc^{(t)}$, and $\pmc^{(t)}\id^{(>N)}\succeq\pmc\,\id^{(>N)}$ turns
\eqref{eq:pmcsetting} at $t$ into \eqref{eq:pmcstatement} at $t$, with no
further hypothesis; at $\Tset=\{\off\}$ the minimum is over one number.

\emph{Step 4 (range and computability).} The range $0\le\pmc^{(t)}\le1$
follows at every setting from the positive part and from an entry at most $1$
in the inner minimum. At a setting governed by Definition~\ref{def:pmc} every
entry is a value of $\Bnd^{(t)}$ or the tail value $T^{(t)}_{m_t}$, and
$\Bnd\le1$ (Theorem~\ref{thm:multiclick-bound}). At a setting governed by
Definition~\ref{def:pmcex} the window entries $\ell^{(t)}(n)$ may exceed $1$,
but the inner minimum of each branch contains an analytic remainder at most
$1$: $\Bnd^{(t)}(m^{\ast}_t+1)\le1$ on branch~1, and
$T^{(t)}_{m_t}=1-\sum_k\beta^{(t)}_k\big(d^{(t)}_k\big)^{m_t}\le1$ on
branch~2 because $\beta^{(t)}_k,d^{(t)}_k\ge0$ (the \emph{Range} paragraph of
Remark~\ref{rem:exacteig}). At a setting governed by
Definition~\ref{def:enclosed-weight} the range is Step~2a. And $\pmc$, a
minimum of finitely many
numbers in $[0,1]$, lies there with them. Computability of each $\pmc^{(t)}$
is the finiteness of its window, and of $\pmc$ the finiteness of $\Tset$.
\end{proof}
